\section{Virtualização Baseada em Modo Máquina}

O \textit{Machine Mode} foi concebido originalmente na arquitetura RISC-V como o nível de privilégio mais elevado, responsável pela execução do firmware base e pelo controle direto dos recursos de hardware.
O software residente nesse nível encarrega-se da inicialização da plataforma, da configuração dos subsistemas fundamentais e, subsequentemente, da transferência de controle para o \gls{so} hospedeiro. 

Além das responsabilidades de inicialização, determinadas implementações da arquitetura requerem que o modo máquina forneça serviços privilegiados às camadas superiores. Esses serviços são disponibilizados
por meio de interfaces padronizadas, como a \gls{sbi}, permitindo que sistemas executados em níveis menos privilegiados acessem de forma controlada as funcionalidades de hardware restritas.

No escopo deste trabalho, o M-mode é empregado de forma atípica: como a camada base para a implementação de um hipervisor, visando contornar a ausência da extensão de virtualização (H) na plataforma utilizada.
Esta estratégia fundamenta-se na autoridade irrestrita que o modo máquina possui sobre os demais níveis de privilégio da arquitetura, o que viabiliza o controle e a orquestração integral dos \glspl{so} convidados.
Para tanto, o M-mode dispõe de um conjunto isolado de \glspl{csr} e de mecanismos próprios para o tratamento de interrupções e temporização, operando de maneira completamente independente dos modos que operam acima dele.

O isolamento e o controle dos sistemas convidados baseiam-se fortemente no mecanismo de tratamento de exceções da arquitetura. Sempre que o \gls{so} virtualizado executa uma instrução privilegiada não permitida
ou tenta acessar um recurso restrito, o hardware gera uma exceção que é direcionada nativamente ao modo máquina. Esse comportamento viabiliza a técnica clássica de \textit{trap-and-emulate}, permitindo que o hipervisor
intercepte a operação falha e emule o seu comportamento de maneira controlada.

Embora o M-mode careça dos mecanismos de tradução de memória em dois estágios - como a \gls{mmu} de segundo estágio fornecida pela extensão H -, o isolamento de memória entre os convidados é viabilizado através 
dos dispositivos de \gls{pmp}. Configurado exclusivamente pelo modo máquina, o \gls{pmp} permite delimitar e proteger regiões específicas da memória física contra acessos
indevidos de camadas menos privilegiadas.

Essa granularidade do \gls{pmp} é também empregada na emulação de \gls{mmio}, recurso indispensável para o suporte a periféricos virtuais como o \textit{VirtIO}. Esse mecanismo opera mediante a
configuração das regiões correspondentes aos dispositivos virtuais como inacessíveis ao sistema convidado; consequentemente, qualquer tentativa de leitura ou escrita engatilha uma exceção de falha de acesso,
permitindo que o hipervisor intercepte a operação e processe a requisição para o \gls{so} convidado.
